
\section{Reconcile the recorded statements before annual completion}
\label{sec:recorded-shared-statements}

The mixed return and lease sequence already determines useful financial
quantities. Its cash balance is 361, its lease liability has been paid, and its
recovered stock has been sold. Waiting for a complete annual tax calculation
should not make those recorded movements disappear. At the same time, a
balanced trial balance cannot establish that all required tax, valuation or
closing adjustments have been posted. The reporting question is therefore
precise: what do the completed events explain, and what remains unassessed?

\path{RecordedSharedBookStatements/1} answers that question for the synthetic
book. It replays the opening and requested event sequence, then constructs
financial views of the completed prefix. The supplied closing balance is never
an input. Each view identifies its last completed date, requested and completed
event counts, any failed event, and the declared fiscal period. Even when every
requested event succeeds, complete fiscal coverage and annual tax remain
unresolved. A short successful list does not prove that nothing else happened.

\subsection{A recorded balance sheet and equity movement}

The opening must start at the declared period boundary, with zero income
accounts and no prior event, cash-flow or settlement-tax-adjustment history in
the current execution. All 25 accounts receive exactly one classification.
The original return classes are extended by the right-of-use asset, lease
liability and lease gains; those additions are not hidden in a residual.
Owner-capital accounts must remain unchanged for these selected methods.

At the end of the five-event sequence, recorded assets total 701: cash 361,
inventory 60, property, plant and equipment 100, investments 80, deferred-tax
assets 10 and right-of-use assets 90. Other asset accounts are zero. Recorded
liabilities total 104, comprising debt 100 and deferred-tax liabilities 4.
The refund and lease liabilities are zero after their respective settlements.
All amounts are USD in the supplied synthetic example.

Recorded profit is revenue and gains less the expenses actually posted. The
sale contributes revenue 30 and the lease reduction a gain of 20. Cost of sales
is 16 after the settlement credit and subsequent sale, selling expense is 8,
interest is 11 and depreciation is 10. Current and deferred tax expense accounts
contain no postings in this sequence. Thus recorded profit is
\[
 P_{\mathrm{recorded}}=(30+20)-(16+8+11+10)=5.
\]
The zero tax postings describe the ledger; they are not a calculation of the
tax charge required for the period. Tax assessment stays unresolved alongside
the recorded result.

Let $A$ and $L$ be the totals of recorded assets and liabilities, and let $E_0$
be opening equity after reversing the debit-positive sign of the equity
accounts. The unchanged owner accounts give $E_0=592$. The two independent
reconciliations then agree:
\[
 A-L=701-104=597,
 \qquad E_0+P_{\mathrm{recorded}}=592+5=597.
\]
The first derives equity from the closing balance sheet; the second explains
its change from the recorded income. Any disagreement refuses the report. No
balancing reserve or owner contribution is inferred to make the equations hold.

\subsection{Explain the same cash movement in two ways}

Gross cash origins preserve the distinction between operating activity and
financing. Operating receipts are 30 from the recovered-stock sale. Operating
payments are 40 of refunds, 8 of handling and 11 of interest, totaling 59.
Lease principal of 110 is a financing payment under the retained historical
interest-classification assumption. There are no investing cash movements.
The direct cash view therefore gives
\[
 \mathrm{CFO}=30-59=-29,\qquad
 \Delta\mathrm{cash}=-29-110=-139=361-500.
\]
The complete gross flow records remain attached to the result, so these totals
can be traced back to the individual events rather than to a net cash plug.

The indirect view starts with recorded profit and removes its noncash effects.
Depreciation of 10 is added back. The gain of 20 from reducing the lease is
subtracted because the reduction produced no cash receipt. Releasing the
recovery asset of 16 adds 16 in the working-capital bridge. Settling the refund
liability reduces that credit balance from 40 to zero and uses 40 of cash.
Inventory returns to its opening amount after the sale, and the other selected
working-capital balances do not change. Consequently,
\[
 \mathrm{CFO}
 =\underbrace{5}_{\text{recorded profit}}
  +\underbrace{10}_{\text{depreciation}}
  -\underbrace{20}_{\text{noncash gain}}
  +\underbrace{16}_{\text{recovery asset released}}
  -\underbrace{40}_{\text{refund liability settled}}
 =-29.
\]
The sign of the refund term follows from the accounting convention. The signed
liability moves from $-40$ to zero; subtracting its change gives $-40$ in the
cash bridge. The credit-negative \path{other_gains} balance similarly removes
the positive noncash gain from profit. The reporting code uses the account
changes and gross flows to compute both reconciliations. The displayed values
are the checked example, not constants used to construct a report.

The inherited bridge also retains recorded deferred-tax expense and any
already-expensed opening interest paid as separate terms. The selected
successor does not execute those additional methods. They remain zero in this
example, and the fresh-opening requirement excludes prior paid interest from
the opening. Full tax measurement requires its own supported method and facts.

\subsection{A failed event leaves a smaller report, not a completed year}

In the alternative opening with cash 130, the return settlement leaves 82 and
the later payment of 121 fails. The recorded views describe only the lease
reduction and return settlement that completed. They show profit 16 and
operating cash outflow 48, with the failed payment and one unexecuted event
still visible. Their last completed date is 1 February 2024. They do not show
the lease as paid or describe the attempted payment date as a completed close.

If the first event is unsupported, a valid fresh opening can still be shown
with zero completed events and its explicit refusal. An invalid opening
supplies no financial views. An empty valid event list supplies the opening
view, but proves neither complete period coverage nor zero annual activity.
These distinctions preserve information without manufacturing a closing state.

All results keep tax assessment, period completion and annual completion
unresolved. The original native annual-report function retains its separate
tax-close prerequisite and still refuses the incomplete-tax control. The new
result is a recorded-prefix report with \path{annual_close_ready} and
\path{complete_financial_statements} false. Its arithmetic can be exact even
when required accounting adjustments remain unknown.

Complete replay checks statements, classifications, execution, gross origins,
source identity and coverage together. A changed profit, dropped flow, omitted
account, forged completion flag or altered component memory fails verification.
The implementation and complete requests/results are retained in
\path{scripts/shared_book_statements_v1.py} and
\path{data/accounting-master/shared-book-statements-20261024}.
This supplies the recorded-statement connection for the finite joint book;
complete annual methods and actual-source admission remain further work.
